\exercisesection{Depth and homological dimension}

\begin{enumerate}
\def\labelenumi{\arabic{enumi})}
\tightlist
\item
  Let \(\rho: A \to B\) be a local homomorphism of Noetherian local rings, \(N\) a finitely generated B-module. Suppose that the A-modules \(B\) and \(N\) are flat.
\end{enumerate}

\begin{enumerate}
\def\labelenumi{\alph{enumi})}
\item
  Prove the equality \(\mathrm{dp}_{\mathrm{B}}(\mathrm{N}) = \mathrm{dp}_{\kappa_{\mathrm{A}} \otimes_{\mathrm{A}} \mathrm{B}}(\kappa_{\mathrm{A}} \otimes_{\mathrm{A}} \mathrm{N})\) (consider a resolution of N by finitely generated free B-modules).
\item
  Let M be a finitely generated A-module. Prove the equality
\end{enumerate}

\[
\mathrm{dp} _ {\mathrm{B}} (\mathrm{M} \otimes_ {\mathrm{A}} \mathrm{N}) = \mathrm{dp} _ {\mathrm{A}} (\mathrm{M}) + \mathrm{dp} _ {\mathrm{B}} (\mathrm{N}).
\]

(Let K (resp. L) be a resolution of the A-module M (resp. of the B-module N) by finitely generated free modules; prove that \(K \otimes_{A} L\) is a resolution of \(M \otimes_{A} N\) by free B-modules. If dp \(_{A}\) (M) = m and dp \(_{B}\) (N) = n, compute \(\mathrm{H}_{m+n}(\kappa_{\mathrm{B}} \otimes_{\mathrm{B}} (\mathrm{K} \otimes_{\mathrm{A}} \mathrm{L}))\) in terms of \(\mathrm{H}_{m}(\kappa_{\mathrm{A}} \otimes_{\mathrm{A}} \mathrm{K})\) and \(\mathrm{H}_{n}(\kappa_{\mathrm{B}} \otimes_{\mathrm{B}} \mathrm{L})\) .

\begin{enumerate}
\def\labelenumi{\arabic{enumi})}
\setcounter{enumi}{1}
\tightlist
\item
  Let A be a Noetherian ring, M an A-module admitting a finite free resolution (L, p) by finitely generated free modules.
\end{enumerate}

\begin{enumerate}
\def\labelenumi{\alph{enumi})}
\tightlist
\item
  Show that for every prime ideal p of A, one has
\end{enumerate}

\[
\sum_ {i \geqslant 0} (- 1) ^ {i} \operatorname{rg} _ {\mathrm{A}} (\mathrm{L} _ {i}) = \sum_ {i \geqslant 0} (- 1) ^ {i} \left[ \operatorname{Tor} _ {i} ^ {\mathrm{A}} (\mathrm{M}, \kappa (\mathfrak {p})): \kappa (\mathfrak {p}) \right] = \sum_ {i \geqslant 0} (- 1) ^ {i} \left[ \operatorname{Ext} _ {\mathrm{A}} ^ {i} (\mathrm{M}, \kappa (\mathfrak {p})): \kappa (\mathfrak {p}) \right];
\]

denote \(\chi(\mathrm{M})\) this integer.

\begin{enumerate}
\def\labelenumi{\alph{enumi})}
\setcounter{enumi}{1}
\item
  Prove that one has \(\chi(M) \geqslant 0\) (consider an ideal \(\mathfrak{p}\) associated to A).
\item
  Prove that the following conditions are equivalent:
\end{enumerate}

\begin{enumerate}
\def\labelenumi{(\roman{enumi})}
\item
  \(\chi (\mathbf{M}) = 0\)
\item
  the annihilator of M is not reduced to zero;
\item
  the annihilator of M contains a cancellable element.
\end{enumerate}

(If \(\chi(M) = 0\) , prove that \(M_{\mathfrak{p}}\) is zero for every \(\mathfrak{p} \in \operatorname{Ass}(A)\) ; if \(\chi(M) > 0\) , prove similarly that the annihilator of Ann(M) contains a cancellable element, which implies Ann(M) = 0.)

\begin{enumerate}
\def\labelenumi{\arabic{enumi})}
\setcounter{enumi}{2}
\tightlist
\item
  Let A be a Macaulay local ring. Recall (A, II, p.186, exerc. 16) that an ideal of A is said to be irreducible in A if it is distinct from A and is not the intersection of two ideals strictly containing it. Prove that the following conditions are equivalent:
\end{enumerate}

\begin{enumerate}
\def\labelenumi{(\roman{enumi})}
\item
  A is a Gorenstein ring;
\item
  every ideal of A generated by a maximal secant sequence is irreducible;
\item
  there exists a maximal secant sequence in A generating an irreducible ideal.
\end{enumerate}

(Reduce to the case where A has dimension 0, and consider the socle \(\mathrm{Hom}_{\mathrm{A}}(\kappa_{\mathrm{A}},\mathrm{A})\) of A.)

\begin{enumerate}
\def\labelenumi{\arabic{enumi})}
\setcounter{enumi}{3}
\tightlist
\item
  Let A be a Noetherian ring, M a nonzero finitely generated A-module. One has grade(M) ≤ dpA(M); M is said to be perfect if equality holds.
\end{enumerate}

\begin{enumerate}
\def\labelenumi{\alph{enumi})}
\tightlist
\item
  Let d be an integer; prove that the following conditions are equivalent:
\end{enumerate}

\begin{enumerate}
\def\labelenumi{(\roman{enumi})}
\item
  the A-module M is perfect of projective dimension d;
\item
  we have \(\operatorname{Ext}_{\mathrm{A}}^{p}(\mathrm{M},\mathrm{A}) = 0\) for \(p\neq d\)
\item
  for every prime (resp. maximal) ideal \(\mathfrak{p}\) in the support of M, the \(\mathrm{A}_{\mathfrak{p}}\) -module \(\mathrm{M}_{\mathfrak{p}}\) is perfect of projective dimension \(d\) .
\end{enumerate}

(First prove the equivalence of (i) and (ii).)

\begin{enumerate}
\def\labelenumi{\alph{enumi})}
\setcounter{enumi}{1}
\item
  Let J be an ideal of A generated by a completely secant sequence, and let m be an integer \(\geqslant 1\) . Prove that the A-module \(A/J^{m}\) is perfect.
\item
  If M is perfect, its associated prime ideals are the ideals p such that \(\text{prof}(A_{\mathfrak{p}}) = \text{dp}_{\mathrm{A}}(M)\) .
\item
  For every perfect module N of projective dimension d, set \(\mathrm{N}^{\prime} = \mathrm{Ext}_{\mathrm{A}}^{d}(\mathrm{N}, \mathrm{A})\) . Prove that \(N^{\prime}\) is perfect, of projective dimension d, and that \((\mathrm{N}^{\prime})^{\prime}\) is isomorphic to N (consider a projective resolution of N). We have \(\mathrm{Ass}(\mathrm{N}^{\prime}) = \mathrm{Ass}(\mathrm{N})\) .
\item
  Let \(x\) be a non-zero-divisor in the radical of A, such that the homothety \(x_{\mathrm{M}}\) is injective. For the A-module M to be perfect of projective dimension \(d\) , it is necessary and sufficient that the (A/xA)-module M/xM be so.
\item
  A Macaulay module of finite projective dimension is perfect (reduce to the local case, and argue by induction on dim(M)).
\item
  If A is a Macaulay ring, the perfect modules are the Macaulay modules of finite projective dimension.
\end{enumerate}

\begin{enumerate}
\def\labelenumi{\arabic{enumi})}
\setcounter{enumi}{4}
\tightlist
\item
  \begin{enumerate}
  \def\labelenumii{\alph{enumii})}
  \tightlist
  \item
    Let \(k\) be a field, and R a \(k\) -algebra of finite degree. For R to be a Gorenstein ring, it is necessary and sufficient that the R-module \(\mathrm{Hom}_{k}(\mathrm{R}, k)\) be free of rank 1 (reduce to the case where R is local, and observe that the R-module \(\mathrm{Hom}_{k}(\mathrm{R}, k)\) is injective and contains a submodule isomorphic to \(\kappa_{\mathrm{R}}\) ; for another proof, cf. § 8, no. 6, cor. of prop. 5).
  \end{enumerate}
\end{enumerate}

\begin{enumerate}
\def\labelenumi{\alph{enumi})}
\setcounter{enumi}{1}
\item
  Assume that the above conditions are satisfied, and moreover that the algebra R is graded of type N, with \(R_{0} = k\) . Let s be the largest integer such that \(R_{s} \neq 0\) . Prove that \(R_{s}\) is of dimension 1 over k, and equal to the socle of R; for every non-zero k-linear form \(\ell\) on \(R_{s}\) , the k-bilinear form \((x, y) \mapsto \ell(xy)\) on \(R_{i} \times R_{s- i}\) induces for every i an isomorphism of \(R_{i}\) on \(\operatorname{Hom}_{k}(R_{s-i}, k)\) .
\item
  Let A be a Noetherian ring, and B an A-algebra which is a flat module of finite type. Prove that the following conditions are equivalent:
\end{enumerate}

\begin{enumerate}
\def\labelenumi{(\roman{enumi})}
\item
  the B-module \(\operatorname{Hom}_{\mathrm{A}}(\mathbf{B},\mathbf{A})\) is projective of rank 1;
\item
  for every maximal (resp. prime) ideal \(\mathfrak{p}\) of A, the ring \(\kappa(\mathfrak{p})\otimes_{\mathrm{A}}\mathrm{B}\) is a Gorenstein ring.
\end{enumerate}

These conditions are satisfied in particular when B is a Gorenstein ring.

¶ 6) Let A be a Noetherian local ring, M a finitely generated A-module of finite projective dimension, and N an A-module such that \(\operatorname{Tor}_{i}^{\mathrm{A}}(\mathrm{M},\mathrm{N}) = 0\) for i \textgreater{} 0.

\begin{enumerate}
\def\labelenumi{\alph{enumi})}
\tightlist
\item
  Define for every integer n a canonical isomorphism
\end{enumerate}

\[
\operatorname{Ext} _ {\mathrm{A}} ^ {n} (\kappa_ {\mathrm{A}}, \mathrm{M} \otimes_ {\mathrm{A}} \mathrm{N}) \longrightarrow \bigoplus_ {p \geqslant n} \operatorname{Tor} _ {p - n} ^ {\mathrm{A}} (\kappa_ {\mathrm{A}}, \mathrm{M}) \otimes_ {k} \operatorname{Ext} _ {\mathrm{A}} ^ {p} (\kappa_ {\mathrm{A}}, \mathrm{N}),
\]

and prove that one has \(\mathrm{di}_{\mathrm{A}}(\mathrm{M}\otimes_{\mathrm{A}}\mathrm{N})\leqslant \mathrm{di}_{\mathrm{A}}(\mathrm{N})\).

(Let (P,p) be a finite free resolution of M, (I,e) an injective resolution of N, and C the complex \(P \otimes_A I\). Prove that the complex

\[
0 \to \mathrm{C} ^ {0} / \mathrm{B} ^ {0} (\mathrm{C}) \to \mathrm{C} ^ {1} \to \mathrm{C} ^ {2} \to \dots
\]

defines an injective resolution of \(M \otimes_A N\).)

\begin{enumerate}
\def\labelenumi{\alph{enumi})}
\setcounter{enumi}{1}
\item
  Prove the inequality \(\mathrm{prof}_{\mathrm{A}}(\mathrm{N}) = \mathrm{dp}_{\mathrm{A}}(\mathrm{M}) + \mathrm{prof}(\mathrm{M} \otimes_{\mathrm{A}} \mathrm{N})\) . Deduce another proof of the Auslander-Buchsbaum theorem.
\item
  Assume moreover that N is of finite projective dimension. Prove the equality \(\mathrm{dp}_{\mathrm{A}}(\mathrm{M}\otimes_{\mathrm{A}}\mathrm{N})=\mathrm{dp}_{\mathrm{A}}(\mathrm{M})+\mathrm{dp}_{\mathrm{A}}(\mathrm{N})\) .
\item
  Assume that A is a Gorenstein ring, and denote its dimension by d. Let n be an integer; show that the \(\kappa_{A}\) -vector spaces \(\operatorname{Ext}_{\mathrm{A}}^{n}(\mathbf{k}_{\mathrm{A}},\mathrm{M})\) , \(\operatorname{Tor}_{d-n}^{\mathrm{A}}(\mathbf{k}_{\mathrm{A}},\mathrm{M})\) and \(\operatorname{Ext}_{\mathrm{A}}^{d-n}(\mathrm{M},\mathbf{k}_{\mathrm{A}})\) have the same dimension.
\end{enumerate}

\begin{enumerate}
\def\labelenumi{\arabic{enumi})}
\setcounter{enumi}{6}
\tightlist
\item
  Let A be a Noetherian local ring, M a finitely generated A-module, and T a finitely generated A-module of finite injective dimension. Prove the equality
\end{enumerate}

\[
\operatorname{prof} (\mathrm{A}) = \operatorname{prof} _ {\mathrm{A}} (\mathrm{M}) + \sup \left\{i \mid \operatorname{Ext} _ {\mathrm{A}} ^ {i} (\mathrm{M}, \mathrm{T}) \neq 0 \right\}
\]

(reason by induction on \(\mathrm{prof}_{\mathrm{A}}(\mathrm{M})\) , using prop. 9).

\begin{enumerate}
\def\labelenumi{\arabic{enumi})}
\setcounter{enumi}{7}
\tightlist
\item
  Let A be a ring, and \(u: \mathrm{A}^p \to \mathrm{A}^q\) an A-linear map. We call the rank of \(u\) , and denote it by \(\text{rg}(u)\) , the largest integer \(r\) such that \(\wedge^r u \neq 0\) . One denotes by \(\mathfrak{d}(u)\) the ideal of A generated by the minors of order \(\text{rg}(u)\) of the matrix of \(u\) .
\end{enumerate}

\begin{enumerate}
\def\labelenumi{\alph{enumi})}
\item
  Show that for \(\mathfrak{d}(u)\) to be equal to A, it is necessary and sufficient that the cokernel of \(u\) be projective of rank \(q - \mathrm{rg}(u)\) (reduce to the case where \(\mathrm{A}\) is local and write the matrix of \(u\) in a suitable basis).
\item
  Let \(v: \mathbf{A}^q \to \mathrm{A}^s\) an A-linear map; we assume that we have \(\mathfrak{d}(u) = \mathfrak{d}(v) = \mathbf{A}\) For \(\operatorname{Ker} v = \operatorname{Im} u\) , it is necessary and sufficient that one has \(\operatorname{rg}(u) + \operatorname{rg}(v) = q\) (same method).
\item
  If \(\operatorname{rg}(u) = q\) , prove that \(\mathfrak{d}(u)\) annihilates the cokernel of \(u\) (by reducing to the case \(p = q\) ).
\end{enumerate}

Let A be a ring, and let (L, d) be a bounded complex of finitely generated free A-modules, zero in degrees \textless{} 0. We propose to prove that the following conditions are equivalent:

\begin{enumerate}
\def\labelenumi{(\roman{enumi})}
\item
  We have \(\mathrm{H}_i(\mathrm{L}) = 0\) for \(i > 0\) ;
\item
  We have \(\operatorname{prof}_{\mathrm{A}}(\mathfrak{d}(d_i); \mathrm{A}) \geqslant i\) and \(\operatorname{rg}(d_i) + \operatorname{rg}(d_{i+1}) = \operatorname{rg}_{\mathrm{A}}(\mathrm{L}_i)\) for every \(i \geqslant 1\) .
\end{enumerate}

\begin{enumerate}
\def\labelenumi{\alph{enumi})}
\item
  Show that it suffices to prove this statement when the ring A is Noetherian (consider the subring of A generated by the coefficients of the matrices \(d_{i}\) for \(i \geqslant 1\)). From now on, this condition is assumed to hold.
\item
  If A is a local ring of dimension 0, prove that (ii) implies (i) (use Exercise 8, b)).
\item
  Assume that the ring A is local, that condition (ii) is satisfied, and that \(\operatorname{Supp}(\mathrm{H}_i(\mathrm{L})) = \{\mathfrak{m}_{\mathrm{A}}\}\) for \(i \geqslant 1\) . Prove that \(\mathrm{H}_i(\mathrm{L}) = 0\) for \(i \geqslant 0\) (let \(p\) be the largest integer such that \(\mathfrak{d}(d_p) \neq \mathrm{A}\); deduce from exercise 8, b) that \(\mathrm{H}_i(\mathrm{L}) = 0\) for \(i > p\) , then from exerc. 8, a) and from cor. 2 of prop. 3 \(n^{\circ} 2\) ) that the complex \(0 \to L_p / B_p(L) \to L_{p-1} \to \ldots \to L_0\) is acyclic in degrees \(>0\) ).
\item
  Prove the implication (ii)⇒(i) (reduce to the case where A is local, and argue by induction on dim(A)).
\item
  We assume that the complex L is acyclic in degrees \textgreater{} 0. Let p be a prime ideal associated with A. Show that for every \(i \geqslant 1\) (Coker \(d_{i}\) )p is a free Ap-module, of rank \(\sum_{p \geqslant i-1} (-1)^{p} \operatorname{rg}_{\mathrm{A}}(\mathrm{L}_{p})\) ; moreover, we have \(\operatorname{rg}((d_{i})_{\mathfrak{p}}) = \sum_{p \geqslant i} (-1)^{p-i} \operatorname{rg}_{\mathrm{A}}(\mathrm{L}_{p})\) and \(\mathfrak{d}((d_{i})_{\mathfrak{p}}) = A_{\mathfrak{p}}\) (exerc. 8).
\end{enumerate}

\(f)\) Hence we have \(\operatorname{rg}(d_i) = \sum_{p \geqslant i} (-1)^{p - i} \operatorname{rg}_A(L_p)\) and consequently \(\operatorname{rg}(d_i) + \operatorname{rg}(d_{i+1}) = \operatorname{rg}_\mathrm{A}(L_i)\) for every \(i \geqslant 1\) .

\(g)\) Let \(i\) be an integer \(\geqslant 1\) , and \(\mathfrak{p}\) a prime ideal of \(\mathrm{A}\). Show that the condition \(\mathrm{dp}_{\mathrm{A}_{\mathfrak{p}}}(H_0(L)_{\mathfrak{p}}) \geqslant i\) is equivalent to \(\mathfrak{d}(d_i)A_{\mathfrak{p}} \neq A_{\mathfrak{p}}\) that is, \(\mathfrak{d}(d_i) \subset \mathfrak{p}\) (use exerc. 8, a)). Using prop. 8 of § 1, no. 5 and th. 1, deduce that one has \(\operatorname{prof}_{\mathrm{A}}(\mathfrak{d}(d_i); A) \geqslant i\) which proves the implication (i) \(\Rightarrow\) (ii).

¶ 10) We retain the hypotheses of Exerc. 9; we assume that the complex L is exact in degrees ≠ 0.

\begin{enumerate}
\def\labelenumi{\alph{enumi})}
\item
  Prove the inclusion \(V(\mathfrak{d}(d_{i+1})) \subset V(\mathfrak{d}(d_i))\) for every \(i \geqslant 1\) (observe that by Exercise 8, a prime ideal p does not contain \(\mathfrak{d}(d_i)\) if and only if the \(A_{\mathfrak{p}}\)-module \((\operatorname{Coker} d_i)_{\mathfrak{p}}\) is free).
\item
  We now assume that the annihilator of \(\mathrm{H}_{0}(\mathrm{L})\) is not reduced to zero. Prove that one has \(\mathrm{rg}(d_{1}) = \mathrm{rg}_{\mathrm{A}}(\mathrm{L}_{0})\) (if f is a cancellable element of \(\mathfrak{o}(d_{1})\) , the \(A_{f}\) -module \(\mathrm{H}_{0}(\mathrm{L})_{f}\) is projective of constant rank according to exerc. 8, hence necessarily zero); deduce from this that the support of \(\mathrm{H}_{0}(\mathrm{L})\) is \(\mathrm{V}(\mathfrak{o}(d_{1}))\) .
\item
  We set \(p = \text{prof}_A(\mathfrak{d}(d_1); A)\) . Prove that we have \(V(\mathfrak{d}(d_i)) = \text{Supp}(\mathbf{H}_0(L))\) for \(1 \leqslant i \leqslant p\) (reason as in the proof of \(a\) ), observing that if a prime ideal \(\mathfrak{p}\) of the support of \(\mathrm{H}_0(L)\) did not contain \(\mathfrak{d}(d_k)\) , one would have \(\text{dp}_{\mathrm{A}_{\mathfrak{p}}} \, H_0(L)_{\mathfrak{p}} < k \leqslant \text{grade}_{\mathrm{A}_{\mathfrak{p}}} \, H_0(L)_{\mathfrak{p}}\) .
\end{enumerate}

\begin{enumerate}
\def\labelenumi{\arabic{enumi})}
\setcounter{enumi}{10}
\tightlist
\item
  Let A be a local ring.
\end{enumerate}

\begin{enumerate}
\def\labelenumi{\alph{enumi})}
\tightlist
\item
  Let n be an integer \(\geqslant 1\) , \(u : A^{n-1} \to A^n\) an A-linear map, \(\mathfrak{d}(u)\) the ideal generated by the minors of order n-1 of u (exerc. 8). If \(\operatorname{prof}_{\mathrm{A}}(\mathfrak{d}(u); A) \geqslant 2\) , prove that \(\mathfrak{d}(u)\) admits the resolution
\end{enumerate}

\[
0 \rightarrow \mathrm{A} ^ {n - 1} \xrightarrow {u} \mathrm{A} ^ {n} \xrightarrow {v} \mathfrak {d} (u) \rightarrow 0,
\]

where \(v\) is induced by the homomorphism \(\wedge^{n-1}(^t u): A^n \to A\) (apply Exercise 9).

\begin{enumerate}
\def\labelenumi{\alph{enumi})}
\setcounter{enumi}{1}
\tightlist
\item
  Conversely, let J be an ideal of A of projective dimension 1. Prove that there exists a cancellable element a of A, an integer n, and an A-linear map \(u: A^{n-1} \to A^n\) such that \(\mathrm{J} = a\mathfrak{d}(u)\) ; one has \(\operatorname{prof}_{\mathrm{A}}(\mathrm{J}; \mathrm{A}) = 2\) ("Hilbert-Burch theorem": deduce from exerc. 9 that a minimal resolution of J is of the form \(0 \to A^{n-1} \xrightarrow{u} \mathrm{A}^n \to J\) , with \(\operatorname{prof}_{\mathrm{A}}(\mathfrak{d}(u); \mathrm{A}) \geqslant 2\) then apply a) and exerc. 1 of § 1, observing that \(\mathfrak{d}(u)\) contains a cancellable element of A).
\end{enumerate}

\begin{enumerate}
\def\labelenumi{\arabic{enumi})}
\setcounter{enumi}{11}
\tightlist
\item
  Let A be a Noetherian local ring, and let M and N be finitely generated non-zero \(\mathrm{A}\)-modules. Assume the projective dimension of M is finite. Let \(q\) be the largest integer such that \(\operatorname{Tor}_q^A(M, N) \neq 0\) .
\end{enumerate}

\begin{enumerate}
\def\labelenumi{\alph{enumi})}
\item
  Assume \(\mathrm{prof}_{\mathrm{A}}(\mathrm{Tor}_{q}^{\mathrm{A}}(\mathrm{M},\mathrm{N})) = 0\) . Prove the equality \(\mathrm{prof}_{\mathrm{A}}(\mathrm{N}) + q = \mathrm{dp}_{\mathrm{A}}(\mathrm{M})\) (one may reason by induction on \(\mathrm{prof}_{\mathrm{A}}(\mathrm{N})\) , considering an element \(x\) of \(\mathfrak{m}_{\mathrm{A}}\) a non-zero-divisor in N).
\item
  We further assume that the A-module \(M \otimes_{A} N\) is of finite length. Deduce from a) that \(\operatorname{Tor}_{i}^{A}(M, N)\) is zero for i \textgreater{} depth(A) - depth(M) - depth(N), and nonzero when equality holds. In particular, \(\operatorname{prof}(M) + \operatorname{prof}(N) \leqslant \operatorname{prof}(A)\) .
\end{enumerate}

¶ 13) Let A be a Noetherian local ring. Suppose the following condition is satisfied \(^{1}\) :

(GM) for every Noetherian local A-algebra B, there exists a B-module M such that \(\mathfrak{m}_{\mathrm{B}}\mathrm{M}\neq\mathrm{M}\) and \(\operatorname{prof}_{\mathrm{B}}(\mathrm{M})=\dim(\mathrm{B})\) .

\begin{enumerate}
\def\labelenumi{\alph{enumi})}
\item
  Let P be a complex of flat A-modules of finite length \(\ell\) such that \(\operatorname{Supp}(\mathrm{H}(\mathrm{P})) = \{\mathfrak{m}_{\mathrm{A}}\}\) . Prove that we have \(\ell \geqslant \dim(\mathrm{A})\) (use Exerc. 6 of § 1) \(^2\) .
\item
  Let B be a Noetherian ring, \(\rho: A \to B\) a homomorphism, \(^a\rho: \operatorname{Spec}(B) \to \operatorname{Spec}(A)\) the associated map, M a finitely generated A-module. Prove that the codimension of \(^a\rho^{-1}(\operatorname{Supp}(M))\) in \(\operatorname{Spec}(B)\) is \(\leqslant \mathrm{dp}_A(M)\) (applying \(a\) ) to the complex \(L \otimes_A B_q\) , where \(q\) is the minimal prime ideal of a component of \(^a\rho^{-1}(\operatorname{Supp}(M))\) and L a free resolution of M).
\item
  Let M, N be finitely generated A-modules, with \(M \neq 0\) . Prove the inequality
\end{enumerate}

\[
\dim (N) \leqslant \mathrm{dp} _ {A} (M) + \dim (M \otimes_ {A} N)
\]

(intersection theorem: reason by induction on the integer \(d = \dim(M \otimes_{\mathrm{A}} N)\) . If d = 0, apply b) to the ring B = A/Ann(M); in the general case consider an element of A that is regular for N and for M \(\otimes_{A} N\) ).

\begin{enumerate}
\def\labelenumi{\alph{enumi})}
\setcounter{enumi}{3}
\tightlist
\item
  For there to exist an A-module of finite length and finite projective dimension, it is necessary and sufficient that A be a Macaulay ring.
\end{enumerate}

\begin{enumerate}
\def\labelenumi{\arabic{enumi})}
\setcounter{enumi}{13}
\tightlist
\item
  Let \(\mathrm{A}\) be a Noetherian local ring satisfying condition (GM) (exerc. 13), M a finitely generated A-module of finite projective dimension.
\end{enumerate}

\begin{enumerate}
\def\labelenumi{\alph{enumi})}
\tightlist
\item
  Prove the inequalities
\end{enumerate}

\[
\dim (A) \leqslant \mathrm{dp} _ {A} (M) + \dim (M) \quad \dim (A) - \operatorname{prof} (A) \leqslant \dim (M) - \operatorname{prof} (M).
\]

If moreover the module M is perfect (exerc. 4), one has \(\dim(A) = \mathrm{dp}_{A}(M) + \dim(M)\) (use exerc. 15 of § 1).

\begin{enumerate}
\def\labelenumi{\alph{enumi})}
\setcounter{enumi}{1}
\item
  Let \(\mathfrak{p} \in \mathrm{Supp}(M)\) ; if the \(\mathrm{A}_{\mathfrak{p}}\) -module \(\mathrm{M}_{\mathfrak{p}}\) is Macaulay, the ring \(\mathrm{A}_{\mathfrak{p}}\) is a Macaulay ring. This is in particular the case if \(\mathfrak{p}\) is a minimal prime ideal of \(\mathrm{Supp}(M)\) .
\item
  Let \(\mathbf{x} = (x_{1},\ldots ,x_{r})\) be a regular sequence of elements of \(\mathfrak{m}_{\mathrm{A}}\) , which generates an ideal J of finite projective dimension. Prove that the sequence \(\mathbf{x}\) is completely regular for A (let \(\mathfrak{p}\in V(J)\) and let \(\mathfrak{q}\subset \mathfrak{p}\) be a minimal prime ideal of \(V(J)\) ; observe that one has \(\mathrm{prof}(A_{\mathfrak{p}})\geqslant \mathrm{dp}_{A_{\mathfrak{p}}}(A_{\mathfrak{p}} / J_{\mathfrak{p}})\geqslant \mathrm{dp}_{A_{\mathfrak{q}}}(A_{\mathfrak{q}} / J_{\mathfrak{q}})\) and \(\mathrm{dp}_{A_{\mathfrak{q}}}(A_{\mathfrak{q}} / J_{\mathfrak{q}})\geqslant r\) according to \(b\) ), whence finally \(\mathrm{prof}_{\mathrm{A}}(\mathrm{J};\mathrm{A})\geqslant r\) ; conclude that \(\mathbf{x}\) is completely regular with the aid of cor. 1 of th. 1 of § 1, no. 3).
\item
  Prove that every M-regular sequence is A-regular. (Prove that every element a of A whose annihilator is nonzero is not M-regular; for this, reduce by localization to the case where Supp(M) ∩ Supp(a) = \(\{\mathfrak{m}_{A}\}\) ; by applying exerc. 13 a) to the complex \(L \otimes_A A/p\), where L is a free resolution of M and p is an associated prime ideal of a, prove the inequalities \(\operatorname{prof}(A) \leqslant \dim(A/p) \leqslant \operatorname{dp}_{A}(M)\) and deduce from this that prof(M) = 0.)
\item
  For A to be an integral domain, it is necessary and sufficient that it contain a prime ideal of projective dimension \(< +\infty\) (if \(\mathrm{dp}_{\mathrm{A}}(\mathfrak{p}) < +\infty\) , show that A identifies with a subring of the regular ring \(\mathrm{A}_{\mathfrak{p}}\) ).
\end{enumerate}

\begin{enumerate}
\def\labelenumi{\arabic{enumi})}
\setcounter{enumi}{14}
\tightlist
\item
  Let A be a Noetherian local ring, M a finitely generated A-module of finite injective dimension.
\end{enumerate}

\begin{enumerate}
\def\labelenumi{\alph{enumi})}
\item
  Let \(\mathfrak{p} \in \operatorname{Supp}(M)\) . Prove the equality \(\operatorname{prof}(A) = \operatorname{prof}(A_{\mathfrak{p}}) + \dim(A/p)\) . (Let \(p = \operatorname{prof}(A_{\mathfrak{p}})\) and \(q = \dim(A/p)\) ; observe that one has \(p = \operatorname{di}_{A_{\mathfrak{p}}} (M_{\mathfrak{p}})\) , whence \(\operatorname{Ext}_{A_{\mathfrak{p}}}^{p}(\kappa(\mathfrak{p}), M_{\mathfrak{p}}) \neq 0\) and consequently \(\operatorname{Ext}_{A}^{p+q}(\kappa_A, M) \neq 0\) by lemma 3 of § 1, no. 7, which implies \(\operatorname{prof}(A) = \operatorname{di}_A(M) = p + q\) .) If \(\operatorname{Supp}(M) = \operatorname{Spec}(A)\), A is a Macaulay ring.
\item
  Prove the equality grade(M) + dim(M) = prof(A) (use a) and exercise 15 of § 1).
\end{enumerate}

¶ 16) Let A be a Macaulay ring and M an A-module of finite projective dimension. Prove that for M to be flat, it is necessary and sufficient that for every ideal J of A generated by a completely secant sequence, the canonical homomorphism \(J \otimes_{A} M \to JM\) is bijective (prove by induction on the length of the sequence that this condition implies \(\operatorname{Tor}_{i}^{A}(A/J, M) = 0\) for all i \textgreater{} 0, then, by descending induction on n, \(\operatorname{Tor}_{n}^{A}(N, M) = 0\) for every finitely generated A-module N and every integer n \textgreater{} 0). Example: a torsion-free module over a Dedekind ring is flat (cf. VII, § 2, exerc. 14).
% semantic-fragments: legacy-input-seam
\relax{}
